% !TEX TS-program = xelatex
% !TEX encoding = UTF-8 Unicode
% ================================================
%  Weekly Research Report (English version)
%  Compile: xelatex -> bibtex -> xelatex -> xelatex
%
%  NOTE: this version uses a plain `article` base rather than
%  `sjtuweekly`, since that class's cover page / running headers are
%  built for Chinese labels. If you have (or can get) an English
%  variant of sjtuweekly.cls, swap the \documentclass line below and
%  drop the manual title block in favor of \makeweeklytitle.
% ================================================

\documentclass[11pt]{article}
\usepackage[margin=1in]{geometry}
\usepackage{fontspec}
\usepackage{booktabs}
\usepackage{tabularx}
\usepackage{array}
\usepackage{titlesec}
\usepackage{setspace}

% ---- Typography: standard Western academic-memo pairing ----
\setmainfont{Times New Roman}       % body text
% \setsansfont{Arial}                % uncomment if headings should be sans-serif
\onehalfspacing
\titleformat{\section}{\normalfont\bfseries\large}{\thesection}{1em}{}

\renewcommand{\thefigure}{\arabic{figure}}
\renewcommand{\thetable}{\arabic{table}}

\begin{document}

% ---- Title block (manual, since not using sjtuweekly's cover page) ----
\begin{center}
{\Large\bfseries Weekly Research Report}\\[0.5em]
{\large Huo Jinlong}\\[0.2em]
% Week \#\,xx
\end{center}
\vspace{1.5em}

% ================================================
% Sections 1-3 (motivation / approach / expected outcomes) usually
% stay stable across a project and don't need rewriting every week --
% only Sections 4-7 (progress / discussion / questions / plan) need
% weekly updates.
% ================================================

\section{Motivation}

MoE (Mixture-of-Experts) models rely on heavy All-to-All communication
for expert routing during both training and inference. Existing EPS
(electrical packet switching) interconnects suffer from low bandwidth
utilization and congestion under this sparse, dynamic traffic pattern.
OCS (optical circuit switching) offers higher link bandwidth and lower
static latency, but circuit setup and reconfiguration overhead limits
its practical benefit in fine-grained, dynamic routing scenarios. This
work asks: can predicting expert-routing outcomes and pre-triggering
circuit reconfiguration hide OCS's reconfiguration overhead inside the
compute/communication pipeline -- preserving OCS's bandwidth advantage
while avoiding its dynamic-routing weakness?

\section{Approach \& Supporting Work}

The core approach is to drive an OCS circuit simulator with expert
routing traces extracted from real model weights, and evaluate a
``circuit reuse + reconfiguration/compute overlap'' strategy under
realistic load, rather than relying only on synthetic or uniform
routing assumptions.
% Supporting references (add BibTeX keys and cite, e.g. \cite{example_ocs_paper}):
% \cite{example_ocs_paper}
% \cite{example_moe_routing_paper}

\section{Expected Outcomes (Medium/Long-term)}

Over the next 1--2 months:
\begin{enumerate}
  \item Characterize OCS's communication benefit over EPS across batch
        sizes and topology scales, using real-weight-driven OCS simulation.
  \item Determine how routing learnability (e.g., routing stability
        before/after LoRA fine-tuning) affects the realized gain from
        pre-triggering, and turn this into a reproducible evaluation
        methodology.
  \item Combine pre-triggering with vLLM's prefill/decode communication
        patterns to recommend a deployment granularity for real
        inference-serving scenarios.
\end{enumerate}

\section{Progress This Week}

\textbf{Combining real model weights with OCS simulation}

Real routing traces extracted from model weights were fed into the OCS
simulator this week, with preliminary results below.

\begin{table}[htbp]
\centering
\small
\begin{tabular}{lccc}
\toprule
\textbf{Metric} & \textbf{EPS} & \textbf{OCS} & \textbf{$\Delta$} \\
\midrule
Communication time    & 23,269 $\mu$s & 17,091 $\mu$s & --26.6\% \\
OCS extra overhead     & --            & 333 $\mu$s    & -- \\
Circuit reuse rate     & --            & 96.7\%        & -- \\
Reconfiguration time    & --            & 150 $\mu$s (3 cold starts) & -- \\
\bottomrule
\end{tabular}
\caption{Synthetic-expert + OCS comparison: 87/90 flows go optical
after circuit setup, a clear latency advantage over EPS.}
\label{tab:ocs_eps}
\end{table}

\noindent Takeaway: circuit reuse is high; reconfiguration overlaps
with compute under the pipeline schedule, and is fully hidden under
the DBO schedule.

% If there were open problems this week, note them here, e.g.:
% \noindent\textbf{Open issues:} simulator-vs-real-hardware gap,
% specific difficulties in data/weight preprocessing, etc.

\section{Discussion / Reflections}

The circuit reuse rate under real weights (96.7\%) is close to what
was seen in the synthetic setting, suggesting that although real
routing shows some load imbalance, it still exhibits meaningful
temporal locality -- an early, encouraging sign for pre-triggering in
realistic settings. This still needs to be confirmed at larger
topology scales and across more batch configurations.

\section{Questions for Advisor}

\begin{enumerate}
  \item For evaluating routing traces before/after LoRA fine-tuning,
        is there a recommended metric for quantifying routing
        stability?
  \item Any guidance on the range of real model weights (model size,
        number of experts) to use, to keep the benchmark
        representative?
\end{enumerate}

\section{Plan for Next Week}

\begin{enumerate}
  \item Complete a full multi-batch, multi-topology benchmark on the
        Qwen+OCS testbed and recalibrate OCS's benefit under real
        weights.
  \item Compare routing traces before/after LoRA fine-tuning in the
        OCS simulator, to verify how much routing learnability
        actually contributes to OCS pre-triggering.
  \item Discuss the appropriate granularity for OCS pre-triggering in
        inference-serving settings, in light of vLLM's
        prefill/decode communication patterns.
\end{enumerate}

\vspace{1em}
\noindent That covers this week's progress -- happy to take any
feedback or corrections.

% ================================================
% \bibliographystyle{plain}
% \bibliography{bib/database}

\end{document}